% 似地行星（综述）
% license CCBYSA3
% type Wiki

本文根据 CC-BY-SA 协议转载翻译自维基百科\href{https://en.wikipedia.org/wiki/Earth_analog}{相关文章}。

\begin{figure}[ht]
\centering
\includegraphics[width=8cm]{./figures/849251199efbb53d.png}
\caption{地球与金星的演化路径。金星一直是一个典型范例，用来说明一颗外观上类似地球的行星，其演化结果可能会与地球产生何等巨大的差异。} \label{fig_Eaan_1}
\end{figure}
类地行星（Earth analog），也称为“地球双胞胎”或“第二地球”，是指其环境条件与地球相似的行星或卫星。“类地行星”（Earth-like planet）这一术语也常被使用，但该术语有时泛指任何类地行星。

这一可能性对天体生物学家和天文学家而言尤为重要，其基本逻辑在于：一颗行星与地球越相似，它就越有可能具备维持复杂地外生命的能力。因此，这一主题长期以来一直是科学、哲学、科幻文学以及大众文化中的重要猜想和讨论对象。空间殖民以及“太空与人类存续”理念的倡导者，也长期以来一直在寻找可供定居的类地行星。在遥远的未来，人类甚至可能通过行星改造（terraforming）来人工制造一颗类地行星。

在对系外行星进行科学搜索与研究之前，这一可能性主要通过哲学和科幻作品来探讨。哲学家曾提出，由于宇宙尺度极其巨大，某处必然存在一颗与地球几乎完全相同的行星。“平庸原理”认为，像地球这样的行星在宇宙中应当十分常见；而“稀有地球假说”则认为此类行星极其罕见。迄今为止已发现的成千上万个系外行星系统，在整体结构上与太阳系有着显著差异，这在一定程度上支持了稀有地球假说。

2013 年 11 月 4 日，天文学家基于开普勒空间望远镜的数据报告称，在银河系内围绕类太阳恒星和红矮星的宜居带中，可能存在多达 400 亿颗地球大小的行星 $^{[1][2]}$。从统计意义上看，距离地球最近的此类行星可能位于约 12 光年之内 $^{[1][2]}$。2020 年 9 月，天文学家在已确认的 4000 多颗系外行星中，依据天体物理参数以及地球生命的自然演化历史，识别出了 24 颗“超级宜居行星”（即在宜居性方面优于地球的行星）候选者 $^{[3]}$。
\begin{figure}[ht]
\centering
\includegraphics[width=8cm]{./figures/22fbf76c2e490700.png}
\caption{一颗假想的宜居系外行星及其三颗天然卫星的艺术示意图。} \label{fig_Eaan_2}
\end{figure}
自20世纪90年代以来的科学发现，极大地影响了天体生物学领域的研究范围、行星宜居性模型，以及对地外智能生命（SETI）的探索。
\subsection{历史}
\begin{figure}[ht]
\centering
\includegraphics[width=8cm]{./figures/4c512b83e6f4ad53.png}
\caption{珀西瓦尔·洛威尔将火星描绘成一颗干燥但类似地球、并且适合外星文明居住的行星。} \label{fig_Eaan_3}
\end{figure}
\begin{figure}[ht]
\centering
\includegraphics[width=6cm]{./figures/4f1f34a4309a75e0.png}
\caption{地球纳米布沙漠中的沙丘（上），与土卫六上贝莱特地区的沙丘对比} \label{fig_Eaan_4}
\end{figure}
在 1858 年至 1920 年 之间，包括一些科学家在内的许多人认为火星与地球非常相似，只是更加干燥，拥有厚重的大气层、相近的自转轴倾角、轨道和季节，并且存在一个建造了巨大火星运河的火星文明。这些理论由乔瓦尼·夏帕雷利（Giovanni Schiaparelli）、珀西瓦尔·洛厄尔（Percival Lowell） 等人提出。因此，在文学作品中，火星常被描绘成类似于地球沙漠的红色星球。然而，“水手号”（Mariner，1965 年）和“维京号”（Viking，1975–1980 年）空间探测器获得的图像和数据揭示，火星实际上是一个贫瘠、布满撞击坑的世界$^{[4][5][6][7][8][9]}$。尽管如此，随着研究的不断深入，一些将火星与地球类比的观点仍然存在。例如，“火星海洋假说”起源于维京号任务，并在 20 世纪 80 年代得到普及$^{[10]}$。随着过去存在液态水的可能性被提出，生命或许曾在火星上起源，火星也因此再次被视为一个更接近地球的世界。

同样地，直到 20 世纪 60 年代，包括一些科学家在内的许多人仍认为金星是一个更温暖版本的地球，拥有厚重的大气层，可能是炎热而多尘，或是潮湿并覆盖着水云和海洋$^{[11]}$。在科幻作品中，金星常被描绘成与地球相似的世界，许多人还曾推测存在金星文明。然而，这些看法在 20 世纪 60 年代被推翻：最早的空间探测器获取的精确科学数据表明，金星是一个极端炽热的星球，其表面温度约为 462 °C（864 °F）$^{[12]}$，并被酸性大气所包围，表面气压高达 9.2 MPa（1330 psi）$^{[12]}$。

自 2004 年起，“卡西尼–惠更斯号”（Cassini–Huygens）任务逐步揭示，土星的卫星泰坦是宜居带之外最类似地球的天体之一。尽管其化学组成与地球截然不同，但 2007 年对泰坦湖泊、河流以及流水地貌过程的确认，使其与地球的类比研究取得重大进展$^{[13][14]}$。后续观测（包括气象现象）进一步加深了人们对类地行星上可能发生的地质过程的理解$^{[15]}$。

开普勒空间望远镜自 2011 年起开始观测宜居带内潜在类地行星的凌日现象$^{[16][17]}$。尽管该技术为行星的发现和确认提供了更有效的手段，但仍无法最终判定这些候选行星在多大程度上类似地球$^{[18]}$。到 2013 年，多颗半径小于1.5 个地球半径、位于恒星宜居带内的开普勒候选行星被确认。直到 2015 年，首颗围绕类太阳恒星运行、尺寸接近地球的候选行星 Kepler-452b才被正式公布$^{[19][20]}$。

2023 年 1 月 11 日，NASA 科学家报告探测到 LHS 475 b——一颗类地系外行星，也是詹姆斯·韦布空间望远镜发现的第一颗系外行星$^{[21]}$。
\subsection{属性与判据}
寻找“地球类比体”的概率主要取决于被认为应当相似的属性，而这些属性在不同研究中差异很大。通常认为，这样的天体应是一颗类地行星，已有多项科学研究以此为目标展开探索。常被隐含但不限于的判据包括：行星大小、表面重力、恒星大小与类型（如类太阳恒星）、轨道距离及其稳定性、自转轴倾角与自转周期、相似的地理环境、海洋、大气与气候条件、强磁层，甚至是否存在类似地球的复杂生命。如果存在复杂生命，陆地上可能覆盖着大片森林；如果存在智慧生命，部分陆地甚至可能被城市所覆盖。然而，其中一些被假定的条件，考虑到地球自身的演化历史，可能并不常见。例如，地球的大气并非始终富含氧气，而富氧大气本身正是光合作用生命出现后的生物标志。此外，月球的形成、存在及其对这些特征的影响（如潮汐力）也可能使寻找地球类比体变得更加困难。

确定地球类比体的过程，往往需要在多个不确定性量化层面之间进行权衡。正如人类学家文森特·伊亚伦蒂（Vincent Ialenti）在其关于类比推理认识论的研究中所指出的那样$^{[22]}$，一些行星科学家“更愿意跨越时间与空间，凭借某种信念将两个迥异的对象联系起来”，而另一些人则对此持更为谨慎的态度$^{[23]}$。
\subsubsection{大小}
\begin{figure}[ht]
\centering
\includegraphics[width=8cm]{./figures/944eb86e695ec56c.png}
\caption{大小对比：Kepler-20e$^{[24]}$ 和 Kepler-20f$^{[25]}$ 与金星和地球的比较} \label{fig_Eaan_5}
\end{figure}
大小通常被认为是一个重要因素，因为与地球大小相近的行星更有可能在本质上属于类地行星，并且能够保留类似地球的大气层。$^{[26]}$

该列表包括质量位于 $0.8$–$1.9$ 个地球质量范围内的行星，低于这一范围的通常被归类为次地球型行星（sub-Earth），而高于这一范围的则被归类为超级地球（super-Earth）。此外，仅纳入半径已知落在 $0.5$–$2.0$ 个地球半径范围内（即地球半径的一半到两倍）的行星。

根据尺寸判据，按已知半径或质量来看，最接近地球的行星质量天体包括：
\begin{figure}[ht]
\centering
\includegraphics[width=14.25cm]{./figures/40bad47d87005cf0.png}
\caption{} \label{fig_Eaan_6}
\end{figure}
翻译“这一对比表明，仅凭大小作为衡量标准是不够的，尤其是在宜居性方面。还必须考虑温度因素，因为金星以及2012年发现的半人马座αB行星、开普勒-20（2011年发现））、COROT-7（2009年发现）以及开普勒-42的三颗行星（均于2011年发现）都极为炽热，而火星、木卫三和土卫六则是寒冷的世界，这也导致了它们表面和大气条件的极大差异。这些天体的质量为太阳系的卫星数量仅占地球卫星数量的一小部分，而系外行星的质量则很难准确测量。然而，发现类地行星——其大小与地球相当——具有重要意义，因为这可能表明类地行星的出现频率和分布情况。
\subsubsection{类地}
\begin{figure}[ht]
\centering
\includegraphics[width=6cm]{./figures/5c2efcf45e09b95e.png}
\caption{像土星卫星泰坦这样的表面（由“惠更斯”号探测器拍摄）在外观上与地球的洪泛平原具有一定的相似性。} \label{fig_Eaan_7}
\end{figure}